\section{Explicit primitive recursive coding and search invariants}
\label{app:coding}

This appendix supplies the numerical details behind the record notation.
It also explains why the displayed finite logical arguments do not hide
stronger induction.

\subsection{Natural quotient and remainder}

Fix $m>0$. Define the pair $(Q_m(a),R_m(a))$ by recursion on $a$:
\[
 (Q_m(0),R_m(0))=(0,0).
\]
If the current remainder is $R$, set
\[
 (Q',R')=
 \begin{cases}
 (Q+1,0),&R+1=m,\\
 (Q,R+1),&R+1<m.
 \end{cases}
\]
For a total definition, the second branch can be the default branch;
the invariant ensures its displayed inequality. Simultaneous recursion is
implemented by pairing. Quantifier-free induction proves
\[
 a=mQ_m(a)+R_m(a),\qquad 0\le R_m(a)<m.
\]
The update proves the identity directly in each case. Uniqueness follows
because two distinct quotients differ by at least one, while two
remainders in $[0,m)$ differ in absolute value by less than $m$.
Set $Q_0(a)=0$ and $R_0(a)=a$ as total defaults. This establishes the
primitive recursive division and remainder operations used in the proof.

Positive-factor cancellation is elementary arithmetic: if $c>0$ and
$ca=cb$, then $a<b$ would imply $cb\ge ca+c>ca$, and $b<a$ would
give the symmetric contradiction. Decidability of natural order leaves
$a=b$. The multiplication-order laws used here follow by quantifier-free
induction from addition and multiplication definitions.

\subsection{A canonical integer code}

Encode a signed integer by
\[
 \operatorname{enc}(z)=
 \begin{cases}
 2z,&z\ge0,\\
 2(-z)-1,&z<0.
 \end{cases}
\]
This display describes the intended meaning. The internal functions are
defined directly on natural codes. For a code $c$, put
\[
 P(c)=\begin{cases}c/2,&c\text{ even},\\0,&c\text{ odd},\end{cases}
 \qquad
 M(c)=\begin{cases}0,&c\text{ even},\\(c+1)/2,&c\text{ odd}.\end{cases}
\]
Thus $c$ represents $P(c)-M(c)$ and its magnitude is
$\operatorname{mag}(c)=(c+1)\mathbin{\mathrm{div}}2$.
At most one of $P(c),M(c)$ is nonzero.

For natural $u,v$, define the code of their difference by
\[
 E(u,v)=
 \begin{cases}
 2(u\monus v),&v\le u,\\
 2(v\monus u)\monus1,&u<v.
 \end{cases}
\]
Then canonical coded addition, negation, and multiplication are
\begin{align*}
 c\oplus d&=E(P(c)+P(d),M(c)+M(d)),\\
 \operatorname{neg}(c)&=E(M(c),P(c)),\\
 c\otimes d&=E(P(c)P(d)+M(c)M(d),\;
                 P(c)M(d)+M(c)P(d)).
\end{align*}
These definitions use only natural arithmetic and decidable cases. Their
interpretation as the desired integer operations follows by expansion of
differences. Sign cases reduce their laws to natural polynomial equations.
The square magnitude $\operatorname{mag}(c)^2$ is the natural norm
contribution of the coded integer.

For division by $m>0$, first compute natural magnitude quotient and
remainder $a=mq+r$, $0\le r<m$. A nonnegative code gives quotient $q$
and remainder $r$. For the code of $-a$, use
\[
 \begin{array}{c|cc}
 &\text{signed quotient}&\text{nonnegative remainder}\\\hline
 r=0&-q&0\\
 r>0&-(q+1)&m-r.
 \end{array}
\]
Both rows satisfy $-a=mQ+R$ and $0\le R<m$. In the exact-division case
this is the usual signed exact quotient. For zero divisor choose quotient
zero and remainder equal to the input signed integer. The valid descent
branches never divide by zero.

All signs in the human proof can therefore be translated into ordinary
natural terms. The reference implementation's signed integers and Lean's
\code{Int} are convenient host representations of these functions, not
an additional sort in the object theory.

\subsection{Tuple coding and simultaneous recursion}

A convenient pairing, matching the inspected ProveIt interface, is
\[
 \pi(u,v)=(u+v)^2+u.
\]
It has primitive recursive projections on its image. Given a code $c$,
let $s$ be the largest $j\le c$ with $j^2\le c$, computed by a bounded
scan; put $u=c\monus s^2$ and $v=s\monus u$.
For $c=\pi(u,v)$, write $s_0=u+v$. Since
\[
 s_0^2\le c=s_0^2+u\le s_0^2+s_0<(s_0+1)^2,
\]
the computed square-root floor is $s_0$, and both projections recover
the original $u,v$. On codes outside the image their values are still
defined, and no inverse law is needed there.

Nested pairing codes a signed four-tuple, a multiplier with a four-tuple,
or a remainder with an accumulated tuple. A single primitive recursion
on the resulting code implements every simultaneous update. The defining
equations for the projections then recover the displayed record equations.
There is no recursor on an infinite stream and no need for a sequence
comprehension principle.

\subsection{Bounded search as a primitive recursive function}
\label{app:folds}

Let $b(i,\vec u)\in\{0,1\}$ be a \emph{fixed} primitive recursive test.
The following recursion returns the first successful index below $B$,
or $B$ if there is none:
\begin{align*}
 s(0,\vec u)&=0,\\
 s(k+1,\vec u)&=
 \begin{cases}
 s(k,\vec u),&s(k,\vec u)<k,\\
 k,&s(k,\vec u)=k\text{ and }b(k,\vec u)=1,\\
 k+1,&\text{otherwise}.
 \end{cases}
\end{align*}
The defining cases are decided by natural tests. Quantifier-free induction
proves the following conjunction, using Boolean folds for its bounded parts:
\begin{enumerate}
\item $s(k,\vec u)\le k$;
\item if $s(k,\vec u)<k$, then $b(s(k,\vec u),\vec u)=1$;
\item every $i<s(k,\vec u)$ has $b(i,\vec u)=0$;
\item if $s(k,\vec u)=k$, every $i<k$ has $b(i,\vec u)=0$.
\end{enumerate}
At a successor stage an existing hit is preserved. If there is none,
the new candidate either becomes the stored hit or extends the failed
prefix by one. These are exactly the three branches above. Thus a known
successful candidate $j<B$ forces $s(B,\vec u)<B$, and a non-failure
output gives a genuine witness.

For the finite conjunction fold, the elementary facts needed are prefix
extension and bounded extraction. Extension is the defining equation
\[
 \operatorname{All}(k+1)=\operatorname{All}(k)b(k).
\]
If the product is $1$, both bit factors are $1$. Conversely, if both are
$1$, the product is $1$. For extraction at a fixed parameter $j$, induct
on the quantifier-free formula
\[
 j<k\ \land\ \operatorname{All}(k)=1\quad\to\quad b(j)=1.
\]
In the successor step, decide whether $j<k$ or $j=k$ and use the old
prefix or the last factor. This supplies the usual elimination rule for
the displayed bounded universal statement without quantified induction.
To build a finite conjunction from already available pointwise proofs,
induct on its bit equation, restricting the index to the chosen total
bound by a quantifier-free guard.

The same ideas justify the finite incidence counts in~\eqref{eq:count}.
For an explicitly coded injective list, a prefix contains at most one
entry equal to a fixed residue. The prefix-count invariant is a natural
inequality; absence of repeated equalities is a bounded Boolean condition.
Adding the next equality-test bit preserves the inequality, because a
previous hit and a new hit would be a repeated entry. Finite double-sum
interchange follows by induction on the row count, and within the step
by induction on the column count using associativity of addition.

Here $b$ is a metamathematical schema for each particular test actually
used in the finite program. The object theory has the corresponding
fixed function symbols. Nothing in the proof assumes a primitive
recursive universal evaluator for arbitrary primitive recursive program
codes. Such an assumption is unnecessary and would be an invalid shortcut.

\subsection{A dependency account for the final derivation}

The proof can be assembled in the following finite sequence:
\begin{enumerate}
\item Derive natural arithmetic, coded integer operations, pairing,
finite sums, and bounded-search specifications.
\item Derive the Euler polynomial identity, bounded least-divisor facts,
and the bounded B\'ezout lemma.
\item Prove injectivity of the two residue lists, the finite incidence
contradiction, and the seed's norm and size bounds.
\item Prove the parity-halving identity and the centered odd-descent lemmas,
including positivity and exact quotient divisibility.
\item Apply QF induction to the explicit prime-state and rank invariant
at clock $p$; obtain the prime witness identity.
\item Apply QF induction to the factor-state and rank invariant at clock
$n$; take magnitudes and obtain the four unary witness identities.
\item Introduce the four existential witnesses and universally close the
free input variable using intuitionistic logical rules.
\end{enumerate}

Every mathematical dependency in this list has been proved in the text
or reduced here to its explicit elementary quantifier-free induction.
The corresponding proof objects are the remaining Lean engineering task.
